\begin{table}[H]
\setbox2=\vbox{\begin{minipage}{\linewidth}
\centering\footnotesize
\setbox0=\hbox{
\begin{tabular}{@{}lrrrr@{}}
\toprule
Predictor & Math MAE & Code MAE & QA MAE & Mean MAE \\
\midrule
\shortstack[l]{Power (V53\\delivered)} & 0.292 & 0.393 & 0.679 & 0.455 \\
\shortstack[l]{Per-density\\regression} & 0.425 & 0.458 & 0.437 & 0.440 \\
\shortstack[l]{Median development\\curve} & 0.080 & 0.054 & 0.065 & 0.066 \\
Zero change & 0.218 & 0.145 & 0.326 & 0.230 \\
\bottomrule
\end{tabular}
}\ifdim\wd0>\linewidth\resizebox{\linewidth}{!}{\box0}\else\box0\fi
\caption{The table evaluates pruning repeatability on Pythia-2.8B (final checkpoint; the two cached revisions used carry identical weights, so the six cells per capability are three densities measured twice), $d\in\{0.85,0.75,0.65\}$. Entries give mean absolute prediction errors for signed $\Delta L$ in nats; six cells per capability, with equal state, density and capability weights. All four predictors were frozen before target pruning measurements, using the delivered V53 register and the same available inputs $(N_0,D_0,L_{0,c},d)$, with zero target pruning calibration points. Source-free median and zero ignore source covariates. A2 uses V53's fixed ridge anchor fits and linear interpolation. No refitting, test-based selection or outcome censoring. Dense inputs are measured first; responses use each V6 pruning run's checked dense reference. Rows identify frozen predictors; columns give mean absolute errors for each capability and their mean. V53 identifies the original pruning development register. MAE denotes mean absolute error.}
\label{tab:prune_repeat}
\end{minipage}}
\centering\ifdim\dimexpr\ht2+\dp2\relax>.95\textheight\resizebox*{!}{.95\textheight}{\box2}\else\box2\fi
\end{table}
